\chapter{Conception et architecture du systeme}

\section{Vue d'ensemble de l'architecture}
La conception de \ProjectName{} suit une architecture en couches afin d'assurer separation des responsabilites et maintenabilite :
\begin{enumerate}[leftmargin=1.2cm]
	\item couche \textbf{collecte} (scrapers multi-sites) ;
	\item couche \textbf{traitement} (ETL + normalisation) ;
	\item couche \textbf{donnees} (PostgreSQL) ;
	\item couche \textbf{services} (FastAPI) ;
	\item couche \textbf{presentation} (React dashboard).
\end{enumerate}

\figureplaceholder{Diagramme d'architecture logique en 5 couches}{Architecture globale de \ProjectName}

\section{Conception de la base de donnees}
Le schema relationnel est structure autour de trois familles de tables.

\subsection{Tables de reference}
Les tables \texttt{brands}, \texttt{categories}, \texttt{ranges}, \texttt{products}, \texttt{product\_formats}, \texttt{websites} et \texttt{stores} decrivent le catalogue et le perimetre commercial.

\subsection{Tables operationnelles}
Les tables \texttt{product\_urls}, \texttt{raw\_staging}, \texttt{scraped\_prices} portent les flux de collecte bruts et traces.

\subsection{Tables analytiques}
Les tables \texttt{weekly\_price\_summary} et \texttt{exchange\_rates} alimentent les vues de comparaison et de tendance.

\subsection{Table de securite}
La table \texttt{users} gere les comptes applicatifs, les roles, l'activation des comptes et la date de derniere connexion.

\figureplaceholder{ERD relationnel avec cardinalites et contraintes principales}{Modele relationnel de la plateforme}

\section{Regles d'integrite et qualite des donnees}
Plusieurs regles structurantes ont ete integrees :
\begin{itemize}[leftmargin=1.2cm]
	\item unicite des combinaisons produit (brand, category, range) ;
	\item unicite des variantes (product, format, packaging) ;
	\item unicite des URLs actives par couple (site, produit, magasin) ;
	\item exclusion des prix non valides (prix <= 0) ;
	\item indexation des tables volumineuses pour accelerer les consultations.
\end{itemize}

\section{Conception des services API}
Les routes sont regroupees par domaine fonctionnel.

\begin{longtable}{p{2.7cm}p{2.0cm}p{8.4cm}}
\caption{Synthese des endpoints principaux}\label{tab:endpoints}\\
\toprule
\textbf{Module} & \textbf{Methode} & \textbf{Endpoint principal} \\
\midrule
\endfirsthead
\toprule
\textbf{Module} & \textbf{Methode} & \textbf{Endpoint principal} \\
\midrule
\endhead
Auth & POST & /auth/login (authentification) \\
Auth & GET & /auth/me (profil courant) \\
Prices & GET & /prices/kpis, /prices/filters, /prices/summary, /prices/timeseries \\
Operations & GET & /operations/visibility (supervision scraping) \\
Retail Presence & GET & /retail-presence (matrice de couverture) \\
Admin & GET/POST/PUT/DELETE & /admin/product-formats, /admin/product-urls, /admin/users \\
Admin & POST & /admin/lookups/brands, /admin/lookups/categories, /admin/lookups/ranges \\
\bottomrule
\end{longtable}

\section{Conception securite et gestion des roles}
Le systeme d'authentification suit une logique token \texttt{Bearer} :
\begin{itemize}[leftmargin=1.2cm]
	\item verification du mot de passe hashé (bcrypt) ;
	\item generation d'un token signe (HS256) ;
	\item validation du token sur les routes protegees ;
	\item controle de privilege pour les routes d'administration.
\end{itemize}

Deux profils sont definis :
\begin{itemize}[leftmargin=1.2cm]
	\item \textbf{admin} : acces complet y compris administration et supervision ;
	\item \textbf{user} : acces lecture aux pages analytiques.
\end{itemize}

\section{Conception frontale et experience utilisateur}
L'interface est organisee en deux blocs de navigation :
\begin{itemize}[leftmargin=1.2cm]
	\item \textbf{Analytics} : Prices Explorer, Price Trends, Store Comparison.
	\item \textbf{Admin} : Products \& Catalog, URLs, Retail Presence, Operational Visibility, Users.
\end{itemize}

Cette separation ameliore la lisibilite des usages et limite les erreurs de manipulation pour les profils non administrateurs.

\figureplaceholder{Wireframe de navigation sidebar et routage par role}{Conception UX et parcours utilisateur}

\section{Conclusion}
La conception proposee offre un cadre solide pour relier les besoins metier a une implementation technique claire. Elle facilite les evolutions futures, notamment l'ajout de nouvelles sources, de nouveaux indicateurs et de nouvelles vues analytiques.

\clearpage
